\documentclass[11pt,a4paper]{ctexart}
\usepackage{ipara}

\pagestyle{fancy}
\fancyhf{}
\lhead{\small 平面向量图解讲义：课后作业解析}
\rhead{\small 数学一对一}
\cfoot{\small 第 \thepage\ 页\quad 共 \zpageref{LastPage} 页}
\renewcommand{\headrulewidth}{0.4pt}
\renewcommand{\footrulewidth}{0pt}

\begin{document}

\begin{center}
  {\Large\bfseries 平面向量图解讲义：课后作业解析}\\
  \vspace{0.5em}
  适用：高一/高二同步复习，一轮复习专题
\end{center}

\section{作业 1：数量积与夹角}

\begin{exbox}{题目}
已知 \(|\vec a|=2, |\vec b|=3\)，且 \(\vec a\cdot\vec b=3\)。求 \(\vec a\) 与 \(\vec b\) 的夹角。
\end{exbox}

\teachblock{解答}{
由数量积定义：
\[
\vec a\cdot\vec b
=
|\vec a||\vec b|\cos\theta.
\]
因此：
\[
3=2\times3\cos\theta
\quad\Longrightarrow\quad
\cos\theta=\frac12.
\]
因为向量夹角 \(\theta\in[0,\pi]\)，所以
\[
\boxed{\theta=\frac{\pi}{3}}.
\]
}

\teachblock{易错点}{
不要把数量积 \(\vec a\cdot\vec b\) 与模长乘积 \(|\vec a||\vec b|\) 混为一谈。前者带有角度信息，只有当两向量同向时二者才相等。
}

\section{作业 2：向量线性表示与共线}

\begin{exbox}{题目}
已知 \(\vec a,\vec b\) 不共线，且 \(\vec c=2\vec a-\vec b\)，\(\vec d=\lambda\vec a+3\vec b\)。若 \(\vec c\parallel\vec d\)，求 \(\lambda\)。
\end{exbox}

\teachblock{解答}{
因为 \(\vec c\parallel\vec d\)，所以存在实数 \(k\)，使得
\[
\vec d=k\vec c.
\]
代入：
\[
\lambda\vec a+3\vec b
=
k(2\vec a-\vec b)
=
2k\vec a-k\vec b.
\]
由于 \(\vec a,\vec b\) 不共线，所以它们构成一组基底。对应系数必须相等：
\[
\lambda=2k,
\qquad
3=-k.
\]
由 \(k=-3\)，得
\[
\boxed{\lambda=-6}.
\]
}

\teachblock{方法提醒}{
看到“两个向量分别用同一组不共线基底表示，并且两向量平行”，优先使用“对应系数成比例”或引入比例系数 \(k\)，不要直接去算坐标。
}

\section{作业 3：三点共线的向量判定}

\begin{exbox}{题目}
在 \(\triangle ABC\) 中，点 \(D\) 满足
\[
\overrightarrow{AD}
=
\frac13\overrightarrow{AB}
+
\frac23\overrightarrow{AC}.
\]
证明点 \(D\) 在边 \(BC\) 上。
\end{exbox}

\teachblock{解答}{
将 \(\overrightarrow{AD}\) 改写成从点 \(B\) 出发的向量：
\[
\overrightarrow{BD}
=
\overrightarrow{BA}+\overrightarrow{AD}
=
-\overrightarrow{AB}
+
\frac13\overrightarrow{AB}
+
\frac23\overrightarrow{AC}.
\]
整理得：
\[
\overrightarrow{BD}
=
-rac23\overrightarrow{AB}
+
\frac23\overrightarrow{AC}
=
\frac23(\overrightarrow{AC}-\overrightarrow{AB}).
\]
而
\[
\overrightarrow{BC}
=
\overrightarrow{AC}-\overrightarrow{AB}.
\]
所以
\[
\overrightarrow{BD}
=
\frac23\overrightarrow{BC}.
\]
因此 \(\overrightarrow{BD}\parallel\overrightarrow{BC}\)，且两向量有公共起点 \(B\)，故点 \(B,D,C\) 共线。
又 \(0<\frac23<1\)，所以点 \(D\) 在线段 \(BC\) 上。
}

\teachblock{得分句}{
只写“\(\overrightarrow{BD}=\frac23\overrightarrow{BC}\)，所以共线”通常可以得到主要分数，但完整书写应补一句“二者有公共起点 \(B\)，且比例系数为正并小于 1，所以 \(D\) 在线段 \(BC\) 上”。
}

\newpage

\section{作业 4：中点与重心模型}

\begin{exbox}{题目}
在 \(\triangle ABC\) 中，点 \(M\) 是边 \(BC\) 的中点，点 \(G\) 是三角形重心。用 \(\overrightarrow{AB},\overrightarrow{AC}\) 表示 \(\overrightarrow{AG}\)。
\end{exbox}

\teachblock{解答}{
因为 \(M\) 是 \(BC\) 中点，所以：
\[
\overrightarrow{AM}
=
\frac12(\overrightarrow{AB}+\overrightarrow{AC}).
\]
重心 \(G\) 在线段 \(AM\) 上，且
\[
AG:GM=2:1.
\]
因此
\[
\overrightarrow{AG}
=
\frac23\overrightarrow{AM}
=
\frac23\cdot\frac12(\overrightarrow{AB}+\overrightarrow{AC}).
\]
所以
\[
\boxed{
\overrightarrow{AG}
=
\frac13\overrightarrow{AB}
+
\frac13\overrightarrow{AC}}
\]
}

\teachblock{模型记忆}{
若以三角形顶点 \(A\) 为基准，重心向量公式可直接记为
\[
\overrightarrow{AG}
=
\frac13(\overrightarrow{AB}+\overrightarrow{AC}).
\]
但考试中最好知道来源：中点公式 + 重心分中线为 \(2:1\)。
}

\section{作业 5：向量最值与平方展开}

\begin{exbox}{题目}
已知 \(|\vec a|=1\)，\(|\vec b|=2\)。求 \(|\vec a+\vec b|\) 的取值范围。
\end{exbox}

\teachblock{解答一：三角不等式}{
由向量三角不等式：
\[
\bigl||\vec b|-|\vec a|\bigr|
\le
|\vec a+\vec b|
\le
|\vec a|+|\vec b|.
\]
代入模长：
\[
|2-1|
\le
|\vec a+\vec b|
\le
1+2.
\]
因此
\[
\boxed{1\le|\vec a+\vec b|\le3}.
\]
当 \(\vec a,\vec b\) 反向时取得最小值 1；同向时取得最大值 3。
}

\teachblock{解答二：平方展开}{
设两向量夹角为 \(\theta\)，则
\[
|\vec a+\vec b|^2
=
|\vec a|^2+|\vec b|^2+2\vec a\cdot\vec b.
\]
代入：
\[
|\vec a+\vec b|^2
=
1+4+4\cos\theta
=
5+4\cos\theta.
\]
因为 \(-1\le\cos\theta\le1\)，所以：
\[
1\le|\vec a+\vec b|^2\le9.
\]
从而：
\[
1\le|\vec a+\vec b|\le3.
\]
}

\teachblock{方法比较}{
如果题目只问范围，三角不等式最快；如果后续还要研究角度、数量积或等号条件，平方展开法更容易继续推进。
}

\section{作业 6：向量与几何投影}

\begin{exbox}{题目}
已知 \(|\vec a|=3\)，单位向量 \(\vec e\) 与 \(\vec a\) 的夹角为 \(60^\circ\)。求 \(\vec a\) 在 \(\vec e\) 方向上的数量投影。
\end{exbox}

\teachblock{解答}{
向量 \(\vec a\) 在单位向量 \(\vec e\) 方向上的数量投影为：
\[
\vec a\cdot\vec e
=
|\vec a||\vec e|\cos60^\circ.
\]
因为 \(|\vec e|=1\)，所以：
\[
\vec a\cdot\vec e
=
3\times1\times\frac12
=
\boxed{\frac32}.
\]
}

\teachblock{几何理解}{
数量投影可以理解为“\(\vec a\) 在 \(\vec e\) 方向上的带符号影子长度”。夹角小于 \(90^\circ\) 时投影为正；大于 \(90^\circ\) 时投影为负。
}

\section{复盘清单}

\begin{enumerate}[label=\arabic*. , leftmargin=2em, itemsep=0.6em]
  \item 数量积题先判断：我求的是数值、角度，还是垂直关系？
  \item 看到共线/平行时，优先考虑是否能写成 \(\vec a=\lambda\vec b\)。
  \item 看到三角形中点、重心、分点，优先写位置向量或中点公式。
  \item 看到 \(|\vec a\pm\vec b|\) 最值，优先考虑三角不等式或模平方展开。
  \item 看到投影，优先想到数量积的几何定义。
\end{enumerate}

\end{document}
